\chapter{Truy vấn cơ sở dữ liệu}

\section{Mục tiêu và quy ước}
Các câu truy vấn dưới đây minh họa cách khai thác các quan hệ trong mô hình 20 bảng. Đây là truy vấn mẫu theo lược đồ báo cáo, không khẳng định là câu lệnh đang được ứng dụng gọi nguyên văn. Các cột thời gian được nhóm theo giá trị lưu trong database; khi báo cáo theo ngày Việt Nam cần quy đổi múi giờ phù hợp.

\section{Số máy theo cửa hàng và trạng thái}
\textbf{Mục tiêu:} thống kê số máy đang được ghi nhận ở mỗi cửa hàng.
\begin{lstlisting}[language=SQL]
SELECT s.name AS store_name,
       m.status,
       COUNT(*) AS machine_count
FROM stores AS s
JOIN machines AS m ON m.store_id = s.id
GROUP BY s.id, s.name, m.status
ORDER BY s.name, m.status;
\end{lstlisting}
Khóa ngoại \texttt{machines.store\_id} nối mỗi máy với cửa hàng. Nhóm theo cửa hàng và trạng thái để đếm số máy trong từng nhóm.

\section{Các phiên sử dụng đang hoạt động}
\textbf{Mục tiêu:} xem khách hàng, máy và thời điểm bắt đầu của các phiên chưa kết thúc.
\begin{lstlisting}[language=SQL]
SELECT us.id AS session_id,
       c.full_name AS customer_name,
       m.code AS machine_code,
       us.started_at
FROM usage_sessions AS us
JOIN customers AS c ON c.id = us.customer_id
JOIN machines AS m ON m.id = us.machine_id
WHERE us.status = 'ACTIVE'
ORDER BY us.started_at;
\end{lstlisting}
Truy vấn nối phiên với khách và máy; điều kiện trạng thái chỉ lấy phiên đang hoạt động.

\section{Số phiên và tổng chi phí theo khách hàng}
\textbf{Mục tiêu:} tổng hợp mức sử dụng máy theo khách, kể cả khách chưa có phiên.
\begin{lstlisting}[language=SQL]
SELECT c.id,
       c.full_name,
       COUNT(us.id) AS session_count,
       COALESCE(SUM(us.cost), 0) AS total_usage_cost
FROM customers AS c
LEFT JOIN usage_sessions AS us ON us.customer_id = c.id
GROUP BY c.id, c.full_name
ORDER BY total_usage_cost DESC, c.full_name;
\end{lstlisting}
\texttt{LEFT JOIN} giữ lại khách chưa có phiên; \texttt{COALESCE} hiển thị tổng bằng 0 nếu chưa phát sinh chi phí.

\section{Doanh thu sử dụng máy theo ngày}
\textbf{Mục tiêu:} tổng hợp số hóa đơn đã thanh toán và doanh thu sử dụng máy.
\begin{lstlisting}[language=SQL]
SELECT DATE(issued_at) AS invoice_date,
       COUNT(*) AS invoice_count,
       SUM(usage_cost) AS usage_revenue
FROM invoices
WHERE status = 'PAID'
GROUP BY DATE(issued_at)
ORDER BY invoice_date;
\end{lstlisting}
Chỉ hóa đơn trạng thái \texttt{PAID} được tính. Truy vấn dùng \texttt{usage\_cost} để đo doanh thu phiên máy; nếu hóa đơn bao gồm thêm mặt hàng, cần xác định quy tắc doanh thu riêng.

\section{Top món được gọi nhiều nhất}
\textbf{Mục tiêu:} xếp hạng món theo tổng số lượng trong các đơn chưa bị hủy.
\begin{lstlisting}[language=SQL]
SELECT moi.item_id,
       moi.item_name,
       SUM(moi.quantity) AS quantity_ordered,
       COUNT(DISTINCT mo.id) AS order_count
FROM menu_order_items AS moi
JOIN menu_orders AS mo ON mo.id = moi.order_id
WHERE mo.status <> 'CANCELLED'
GROUP BY moi.item_id, moi.item_name
ORDER BY quantity_ordered DESC, moi.item_name
LIMIT 10;
\end{lstlisting}
Bảng chi tiết giữ tên món tại thời điểm đặt. Vì vậy, kết quả phản ánh tên snapshot trong đơn, kể cả khi tên danh mục hiện tại đã thay đổi.

\section{Doanh số POS theo cửa hàng}
\textbf{Mục tiêu:} tính tổng tiền cuối cùng của các đơn không bị hủy.
\begin{lstlisting}[language=SQL]
SELECT s.name AS store_name,
       COUNT(mo.id) AS order_count,
       SUM(COALESCE(mop.final_amount, mo.total_amount)) AS sales_amount
FROM stores AS s
JOIN menu_orders AS mo ON mo.store_id = s.id
LEFT JOIN menu_order_pricing AS mop ON mop.order_id = mo.id
WHERE mo.status <> 'CANCELLED'
GROUP BY s.id, s.name
ORDER BY sales_amount DESC;
\end{lstlisting}
Nếu đơn có bản ghi tính giá, truy vấn dùng \texttt{final\_amount}; nếu chưa có, dùng tổng tiền ghi trên đơn. Cần thống nhất quy tắc trạng thái thanh toán khi triển khai báo cáo doanh thu thực tế.

\section{Nhân viên được phân công theo cửa hàng}
\textbf{Mục tiêu:} xem danh sách nhân viên đang được phân công cho từng cửa hàng.
\begin{lstlisting}[language=SQL]
SELECT s.name AS store_name,
       e.full_name,
       e.role,
       esa.granted_at
FROM employee_store_assignments AS esa
JOIN employees AS e ON e.id = esa.employee_id
JOIN stores AS s ON s.id = esa.store_id
ORDER BY s.name, e.full_name;
\end{lstlisting}
Bảng liên kết cho phép một nhân viên làm tại nhiều cửa hàng và một cửa hàng có nhiều nhân viên.

\section{Kết luận chương}
Các truy vấn minh họa những thao tác JOIN, GROUP BY, tổng hợp và lọc trên các bảng thuộc phạm vi 20 bảng. Khi dùng với database thật, cần đối chiếu lại tên cột, kiểu dữ liệu, quy tắc trạng thái và dữ liệu hiện có.
